The increasing volume of georeferenced information produced and used by Public Administration has intensified the need for technological solutions capable of organizing, integrating, and providing access to territorial data. In this context, web-based Geographic Information Systems (WebGIS) provide an alternative for expanding access to spatial information and supporting urban planning and management activities. However, the development of such platforms involves challenges related to the integration of different technologies, interoperability among geospatial services, system maintainability, and the provision of interfaces suitable for different devices. Within this scenario, this work presents the development of the ``Sobral em Mapas'' platform, a WebGIS solution designed for the management and interactive visualization of geographic information related to the municipality of Sobral. The platform was structured using a decoupled architecture comprising client, server, and geospatial services. The back-end was developed using the PHP Laravel framework and is responsible for implementing a RESTful API, business rules, administrative data management, authentication mechanisms, and access control. The front-end was developed with React following the Single Page Application (SPA) model, while the cartographic functionalities were implemented using Leaflet and React-Leaflet. Geographic layers are provided through integration with the GeoServer map server using standardized geospatial services, particularly the Web Map Service (WMS). The back-end execution environment was also structured using Docker, contributing to the standardization and reproducibility of the deployment process. As a result, a platform was developed that integrates cartographic visualization, layer management, authentication, user control, and administrative management functionalities, providing a technological foundation for organizing and making geospatial data available within the municipal context.

\keywords{\textit{Web development. Geographic Information Systems. WebGIS. Decoupled architecture. GeoServer.}}
